% capitulo5.tex
\chapter{Conclusiones y recomendaciones}

El presente capítulo sintetiza los hallazgos principales de la comparación entre las NTC-DCEA 2017
y las NTC-DCEA 2023 para el diseño de placas base de acero. Las conclusiones se organizan de
acuerdo con los objetivos de la investigación: identificar los cambios normativos relevantes,
evaluar su efecto en el diseño de placas base y cuantificar su impacto mediante barridos
paramétricos sobre los factores $R_{et}$, $R_{ev}$ y la distancia mínima al borde libre.

% ============================================================
% SECCION 5.1
% ============================================================
\section{Conclusiones generales}

La transición de la NTC-DCEA 2017 a la NTC-DCEA 2023 representa una evolución normativa
significativa para el diseño de conexiones columna--cimentación en estructuras de acero. Sin
embargo, el impacto de esta actualización no se distribuye de manera uniforme en todos los
componentes de la conexión. El análisis desarrollado muestra que las ecuaciones básicas de la placa
base, asociadas al bloque de compresión, la tensión en anclas y el espesor requerido, se mantienen
esencialmente estables entre ambas ediciones. Por ello, el cambio normativo no debe interpretarse
como una modificación profunda del modelo clásico de placa base, sino como una ampliación del
tratamiento del sistema de anclaje y del dado de concreto.

El aporte técnico más importante de la NTC-DCEA 2023 se concentra en el reconocimiento explícito del
refuerzo transversal del pedestal. La incorporación de los términos $R_{et}$ y $R_{ev}$ permite
cuantificar una resistencia que en la edición 2017 no era considerada de manera directa, aun cuando
constructivamente el dado contara con estribos capaces de interceptar el cono o el plano lateral de
falla. Esta diferencia modifica la lectura del diseño: el refuerzo del pedestal deja de ser un
detalle constructivo secundario y pasa a formar parte de la capacidad resistente verificable del
sistema.

La investigación también confirma que la NTC-DCEA 2023 no es simplemente más permisiva que la
edición 2017. Aunque aumenta la resistencia disponible en tensión y cortante cuando existe
confinamiento efectivo, también introduce un requisito geométrico más estricto para anclas de alta
resistencia A325 y A449, al pasar de $6d_o$ a $7d_o$ como distancia mínima proporcional al borde. En
consecuencia, la actualización normativa combina dos efectos simultáneos: reconoce capacidad
adicional en dados bien armados, pero penaliza configuraciones geométricamente compactas que pueden
favorecer fallas frágiles hacia el borde libre.

% ============================================================
% SECCION 5.2
% ============================================================
\section{Conclusiones sobre la comparación normativa}

La comparación normativa permitió identificar que los factores de resistencia para aplastamiento del
concreto y flexión de la placa no presentan cambios sustanciales entre 2017 y 2023. Esto explica
por qué el espesor requerido de placa $t_{req}$ permanece prácticamente igual al aplicar ambas
normas bajo las mismas condiciones de carga, geometría y propiedades mecánicas.

El cambio de mayor alcance conceptual es la reubicación del diseño de placas base desde un apéndice
hacia un capítulo principal dentro de la NTC-DCEA 2023. Esta modificación no solo mejora la
claridad documental, sino que refuerza la importancia del diseño de conexiones de base dentro del
comportamiento global de estructuras de acero. La conexión columna--placa--dado se reconoce como un
sistema crítico de transferencia de fuerzas, especialmente en condiciones sísmicas donde las fallas
frágiles del anclaje pueden comprometer la jerarquía resistente.

También se concluye que la NTC-DCEA 2023 hace más explícita la interacción entre las normas de acero
y concreto. El diseño de la placa base ya no puede evaluarse únicamente desde la resistencia de la
placa de acero; debe verificarse junto con la resistencia del concreto, el anclaje, el refuerzo
transversal y los requisitos geométricos del pedestal.

% ============================================================
% SECCION 5.3
% ============================================================
\section{Conclusiones sobre el análisis paramétrico}

El barrido asociado al factor $R_{et}$ mostró que la resistencia al cono de concreto en tensión
aumenta cuando el plano hipotético de falla intercepta estribos del dado. El incremento depende de
la profundidad efectiva de empotramiento $h_{ef}$ y de la densidad del refuerzo transversal. Por
ello, los mayores beneficios se presentan en dados con estribos cerrados a menor espaciamiento y
con anclas suficientemente embebidas para movilizar un volumen de concreto confinado.

El barrido de distancia mínima al borde evidenció que el cambio de $6d_o$ a $7d_o$ para anclas A325
y A449 puede controlar la geometría del dado. Esta conclusión es especialmente importante para
pedestales compactos, zapatas de lindero o configuraciones donde el espacio disponible para colocar
anclas es limitado. En estos casos, una conexión que cumplía con la NTC-DCEA 2017 puede requerir
redimensionamiento bajo la NTC-DCEA 2023.

El barrido asociado al factor $R_{ev}$ confirmó que la NTC-DCEA 2023 incrementa la resistencia al
cono de concreto por cortante al considerar el aporte del refuerzo transversal interceptado por el
plano lateral de falla. Este efecto es más significativo cuando el espaciamiento de estribos es
reducido, pues aumenta el número de barras capaces de participar en el mecanismo resistente. La
densificación del refuerzo transversal se confirma, por tanto, como una estrategia eficaz para
mejorar el desempeño del dado ante cortante basal.

En conjunto, los tres barridos muestran que los cambios normativos relevantes no actúan de forma
aislada. La profundidad de empotramiento, el diámetro del ancla, el grado del acero, el espaciamiento
de estribos y las dimensiones del dado forman un sistema de variables acopladas. El diseño eficiente
de una placa base no consiste únicamente en aumentar el espesor de la placa o el diámetro de las
anclas, sino en equilibrar la resistencia del acero con la capacidad del concreto y el confinamiento
del pedestal.

% ============================================================
% SECCION 5.4
% ============================================================
\section{Aportaciones de la investigación}

La primera aportación de esta tesis es la sistematización de las diferencias normativas entre la
NTC-DCEA 2017 y la NTC-DCEA 2023 aplicadas al diseño de placas base. Esta comparación permite
separar los cambios meramente editoriales de aquellos que producen efectos cuantificables en el
diseño.

La segunda aportación es la implementación de un análisis paramétrico reproducible, capaz de mostrar
en qué condiciones la actualización normativa modifica la resistencia disponible o la geometría
requerida. Este enfoque evita una interpretación únicamente descriptiva de las normas y permite
medir la sensibilidad del diseño ante variables de uso común en la práctica estructural.

La tercera aportación es la identificación de una lectura técnica balanceada de la NTC-DCEA 2023. La
norma no puede clasificarse de manera absoluta como más conservadora o más permisiva: resulta más
favorable cuando existe confinamiento efectivo por estribos, pero más estricta cuando la distancia al
borde es reducida y se emplean anclas de alta resistencia.

% ============================================================
% SECCION 5.5
% ============================================================
\section{Recomendaciones para el diseño estructural}

Se recomienda que el diseño de placas base bajo la NTC-DCEA 2023 no se limite a verificar el espesor
de la placa. El proyectista debe revisar de manera conjunta el grupo de anclas, el dado de concreto,
el refuerzo transversal y las distancias mínimas al borde. Esta revisión integral es indispensable
para evitar que el sistema falle por mecanismos frágiles antes de que la placa o las anclas alcancen
un comportamiento dúctil.

Para conexiones con anclas de alta resistencia A325 o A449, se recomienda verificar desde etapas
tempranas la distancia mínima al borde libre. Esta comprobación debe realizarse antes de cerrar las
dimensiones del pedestal o de la zapata, ya que el cambio de $6d_o$ a $7d_o$ puede obligar a
modificar la geometría en planta del dado.

En dados sometidos a demandas importantes de tensión o cortante, se recomienda considerar el
refuerzo transversal como una variable de diseño y no solo como un requisito constructivo. La
colocación de estribos cerrados con espaciamiento adecuado puede incrementar la resistencia
disponible mediante $R_{et}$ y $R_{ev}$, además de mejorar la capacidad del pedestal para controlar
la propagación de grietas.

También se recomienda documentar explícitamente en las memorias de cálculo qué versión normativa se
emplea y qué criterios fueron determinantes. En especial, debe indicarse si el diseño quedó
controlado por la resistencia del cono en tensión, la resistencia por cortante o la distancia mínima
al borde, pues cada estado límite conduce a decisiones de diseño distintas.

% ============================================================
% SECCION 5.6
% ============================================================
\section{Recomendaciones para trabajos futuros}

Como línea futura de investigación, se recomienda ampliar el análisis paramétrico a otros tipos de
columna, en particular secciones tubulares y cajón, ya que la NTC-DCEA 2023 incorpora disposiciones
específicas para estas configuraciones. Esto permitiría evaluar si los patrones observados para
perfiles tipo IR se mantienen en geometrías de placa y anclaje diferentes.

También se recomienda complementar el enfoque analítico con modelos numéricos de elementos finitos
que permitan estudiar la interacción placa--grout--concreto--anclas. Un modelo tridimensional
calibrado podría ayudar a evaluar la distribución real de presiones, la flexibilidad de la placa, el
efecto de palanca y la formación progresiva de grietas en el pedestal.

Otra línea útil sería contrastar las expresiones normativas con ensayes experimentales de anclas en
pedestales reforzados. En particular, sería valioso estudiar la compatibilidad entre la deformación
necesaria para movilizar el acero transversal y la apertura de grietas del concreto, especialmente
en el caso del aporte $R_{et}$.

Finalmente, se recomienda desarrollar herramientas de cálculo automatizadas que integren las
verificaciones de placa base, anclaje, refuerzo transversal y distancia al borde. Este tipo de
herramientas facilitaría la aplicación correcta de la NTC-DCEA 2023 y reduciría errores de
interpretación en memorias de diseño.

% ============================================================
% SECCION 5.7
% ============================================================
\section{Cierre}

La comparación realizada permite concluir que la NTC-DCEA 2023 representa una actualización
relevante para el diseño de placas base de acero en México. Su principal contribución no consiste en
cambiar radicalmente el cálculo de la placa, sino en reconocer con mayor precisión la participación
del concreto confinado y del refuerzo transversal en el comportamiento del sistema de anclaje.

El diseño resultante es más coherente con una filosofía de desempeño estructural: favorece el uso
de dados adecuadamente confinados, exige mayor cuidado en la geometría de borde para anclas de alta
resistencia y obliga a entender la placa base como parte de una conexión compuesta. Bajo esta
lectura, la actualización de 2023 no solo modifica ecuaciones aisladas, sino que promueve una forma
más integral y responsable de diseñar la transferencia de fuerzas entre la estructura de acero y la
cimentación.
